% !TEX program = xelatex
% SPDX-License-Identifier: MIT (templates/latex/LICENSE). 이 템플릿 틀만 MIT 이고, 여기에 쓴 글은 쓴 사람의 것입니다.
% 한국어 보고서·학위논문 초안 템플릿: XeLaTeX + kotex, 장·차례·참고문헌(biblatex + biber).
\documentclass[11pt,a4paper,oneside]{report}

\usepackage[hangul]{kotex}  % hangul: 날짜·장 제목 등을 한국어 방식으로
\setmainhangulfont{Noto Serif CJK KR}[AutoFakeSlant=0.15]  % 한글 기울임꼴이 없어 기울여 만든다
\setsanshangulfont{Noto Sans CJK KR}[AutoFakeSlant=0.15]
\setmonohangulfont{Noto Sans Mono CJK KR}

\usepackage[margin=28mm]{geometry}
\usepackage{amsmath,amssymb}
\usepackage{graphicx}
\usepackage{booktabs}
\usepackage[backend=biber,style=numeric,sorting=none]{biblatex}
\addbibresource{refs.bib}
\usepackage[hidelinks]{hyperref}

\AtBeginDocument{%
  \renewcommand{\tablename}{표}%
  \renewcommand{\figurename}{그림}%
  \renewcommand{\contentsname}{차례}%
  \renewcommand{\listtablename}{표 차례}%
  \renewcommand{\listfigurename}{그림 차례}%
  \renewcommand{\bibname}{참고 문헌}%
  \renewcommand{\abstractname}{요약}%
}
\DefineBibliographyStrings{english}{references={참고 문헌}, bibliography={참고 문헌}}

\title{보고서 제목}
\author{홍길동}
\date{\today}

\begin{document}
\maketitle
\tableofcontents

\chapter{서론}\label{ch:intro}
이 장에서는 연구의 배경을 설명한다. 참고 문헌은 \cite{knuth1984} 처럼 인용한다.

\section{연구 목적}
연구의 목적은 한국어 문서를 LaTeX 으로 쓰는 방법을 보이는 것이다.

\chapter{방법}\label{ch:method}
\ref{ch:intro}장에서 말한 목적을 이루기 위해 다음 방법을 쓴다\cite{lamport1994}.

\begin{table}[h]
  \centering
  \caption{비교}\label{tab:compare}
  \begin{tabular}{lc}
    \toprule
    항목 & 값 \\
    \midrule
    가 & 1 \\
    나 & 2 \\
    \bottomrule
  \end{tabular}
\end{table}


\printbibliography[heading=bibintoc]
\end{document}
